%=============================================================================!
% 8.2  WHEN THE NUCLEI ARE CLASSICAL                                          !
%=============================================================================!
\section{When the nuclei may be treated as classical}
\label{sec:pe-classical}

The electrons need quantum mechanics; this book moves the nuclei by
Newton's laws, which \cref{sec:ne-atoms} called an approximation. This
section finds when it is a good one, from a comparison of two energies.

Quantum mechanics changes the motion of a vibration in one way that
matters here: a harmonic oscillator (\cref{sec:os-small}) of angular frequency $\omega$
cannot have any energy it likes, only the values $(n + \frac12)\hbar\omega$, with $n = 0,
1, 2, \dots$ and $\hbar = \SI{0.6582}{\electronvolt\femto\second}$ the
reduced Planck constant\cite{Tuckerman2023}. Its energy therefore comes
in steps of $\hbar\omega$, and even in its lowest state, $n = 0$, it keeps the energy
$\frac12\hbar\omega$, its zero-point energy, which a classical oscillator
lacks. At a temperature $T$, the motion of each vibration carries a
typical energy of order $\kB T$, where $\kB$ is Boltzmann's constant and
$\kB T = \SI{0.02585}{\electronvolt}$ at \SI{300}{\kelvin}
(Chapter~12 makes this precise). When $\hbar\omega \ll \kB T$, a typical
vibration spans many steps, too small to notice, and the classical
description is good. When $\hbar\omega \gg \kB T$, the vibration
sits in its lowest state, and a classical model, which lets it lose its
zero-point energy, describes it wrongly.

With $\omega = 2\pi\nu = 2\pi c\tilde\nu$ for a wavenumber $\tilde\nu$
(\cref{sec:os-molecules}), $\hbar\omega = 2\pi\hbar c\,\tilde\nu$, and
$2\pi\hbar c = \SI{1.2398e-4}{\electronvolt\centi\metre}$. At \SI{300}{\kelvin}:
\begin{itemize}
\item the symmetric O--H stretch of water, \SI{3657}{\per\centi\metre}
   (\cref{tab:os-vibrations}), has $\hbar\omega = \SI{0.4534}{\electronvolt}$,
   17.5 times $\kB T$: thoroughly quantum;
\item the stretch of LiF, \SI{910.34}{\per\centi\metre}
   (\cref{sec:ro-molecules}), has $\hbar\omega =
   \SI{0.1129}{\electronvolt}$, 4.4 times $\kB T$;
\item the model lithium atom rattling in its hollow with a period
   $\mathcal{T} = \SI{170.3}{\femto\second}$ (\cref{sec:os-small}) has $\hbar\omega =
   \hbar\,2\pi/\mathcal{T} = \SI{0.0243}{\electronvolt}$, 0.94 times $\kB
   T$: borderline at room temperature, and 2.8 times $\kB T$ at
   \SI{100}{\kelvin}.
\end{itemize}
The lighter the atom and the stiffer its bonds, the larger $\omega$ and
the stronger the quantum character, so the motion of hydrogen is where
classical nuclei fail first: quantum effects of the nuclei can change
measured properties of water and other hydrogen-bonded systems at room
temperature\cite{MarklandCeriotti2018}. A bond between heavier atoms can
still be quantum, as LiF shows, and the rattle of lithium in its hollow
is borderline. The motions this book measures, the turning and drifting
of whole molecules and the hops of lithium and other ions heavier than
hydrogen, we follow with classical nuclei, accepting that the energy of
the fast stretches is then wrong.
Methods that keep the quantum character of the nuclei exist; they are
beyond the scope of this book.

\begin{keyidea}
A vibration may be treated classically when its quantum step
$\hbar\omega = 2\pi\hbar c\,\tilde\nu$ is small compared with $\kB T$. At room
temperature this holds for slow motions, is borderline for the rattle of
lithium in its hollow, and fails for the stretches of chemical bonds,
most of all bonds to hydrogen.
\end{keyidea}

\begin{selfcheck}
Is the C--H stretch, near \SI{3000}{\per\centi\metre}, classical at
\SI{300}{\kelvin}?
\end{selfcheck}
